\documentclass[12pt]{article}
\usepackage[a4paper, total={6.5in, 9.5in}]{geometry}
\usepackage{graphicx}			
\usepackage{amssymb}
\usepackage{amsmath}			
\usepackage{fancyhdr}
\usepackage{enumerate}       
\usepackage{dirtytalk} 
\usepackage{mathrsfs}
\usepackage[english]{babel}
\usepackage{amsthm}
\usepackage{xcolor}
\usepackage[shortlabels]{enumitem}
\usepackage{titlesec}
\titleformat{\section}{\normalfont\Large\bfseries}{\thesection}{0pt}{\Large}
\setlength{\parindent}{0pt}
\setlength{\parskip}{0.8em}

\newtheorem{theorem}{Theorem}[section]
\newtheorem{corollary}{Corollary}[theorem]
\newtheorem{lemma}[theorem]{Lemma}
\newtheorem{definition}[theorem]{Definition}


\begin{document}
% \renewcommand{\thesection}{Exercise \arabic{section}}


\setcounter{section}{0}
\renewcommand{\thesection}{}


\parindent = 0pt


      {\bf Analysis 1 HW 6 Part 2}
%==================

\section{6.2}
$f$ is continuous on $[a, b]$.\\
$f\geq 0$.\\
$\int_{a}^{b} f(x) \, dx=0$.\\

{\bf To prove:} $f(x)=0$ for all $x\in[a, b]$.

\begin{proof}
    FTSOC, assume $f(p)=q>0$ for some $p\in[a, b]$. Since $f$ is continuous, there exists $\delta>0$ such that $|x-p|<\delta \implies |f(x)-q|< q/2$. Consider a partition $P$ with $x_{k-1}=p-\delta/2$, $x_{k}=p+\delta/2$ for some $k$. 
    $$
    \int_{a}^{b} f \, dx \geq L(P, f)\geq m_{k}\Delta x_{k}\geq \frac{q}{2}\delta
    $$
    This makes it impossible for $\int_{a}^{b} f \, dx$ to be $0$. Thus, $f(x)=0$.
\end{proof}
\pagebreak


\section{6.5}

$f:[a, b]\to \mathbb{R}$.\\
$f$ is bounded.\\

{\bf If $f^{2}\in\mathscr R$ on $[a, b]$, does it follow that $f\in \mathscr R$ on $[a, b]$?}

No. Consider the function which is defined to be $1$ at rational inputs and $-1$ at irrational inputs. Clearly, $f^{2}\in\mathscr R$ since it is a constant function. However, $f\not\in \mathscr R$ since $U(P, f)-L(P, f)=2(b-a)$ for all partitions $P$. 

{\bf If $f^{3}\in\mathscr R$ on $[a, b]$, does it follow that $f\in \mathscr R$ on $[a, b]$?}

Yes. $\sqrt[3]{ x }$ is a continuous function on any interval in $\mathbb{R}$. It follows from Theorem 6.11 that $\sqrt[3]{ f^{3}(x) }=f(x)\in \mathscr R$ on $[a, b]$. 
\pagebreak



\section{6.8}

$f\in\mathscr R$ on $[a, b]$ for every $b>a$ where $a$ is fixed.\\
$\int_{a}^{\infty}f(x)  \, dx\equiv \lim_{ b \to \infty }\int_{a}^{b} f(x) \, dx$, if the limit exists in $\mathbb{R}$.\\


{\bf To prove: }If $f\geq 0$ and is monotone decreasing on $[1, \infty)$, 
$$
\int_{1}^{\infty}f(x)  \, dx \text{ converges} \iff \sum_{n=1}^{\infty}f(n) \text{ converges}.
$$


\begin{proof}

Consider the sequence $(t_{n})$ defined by
$$
t_{n}\equiv \int_{1}^{n} f(x) \, dx.
$$
For any $n$, consider the partition $P=\{ 1, 2, \dots, n \}$. We know that $L(P, f)\leq t_{n}\leq U(P, f)$. Since $f$ is monotone decreasing, we can write
$$
\sum_{k=2}^{n} f(k) \leq t_{n}\leq \sum_{k=1}^{n-1} f(k).
$$
Note that all terms are positive. Now, if $\sum f(n)$ converges, it follows from the comparison test that $t_{n}$ converges. Say it converges to $t$. For every $\epsilon>0$, we can find $N$ such that $n\geq N\implies t-t_{n}<\epsilon$. Note that $\int_{a}^{x} f(x) \, dx$ is a monotone increasing function, since $f\geq 0$. Hence, for $x>N$, $t-\epsilon<\int_{a}^{x} f(x) \, dx<t+\epsilon$. Thus, $\lim_{ x \to \infty }\int_{a}^{x} f(x) \, dx$ exists. 

On the other hand, if $\lim_{ x \to \infty }\int_{a}^{x} f(x) \, dx$ exists, then $t_{n}$ converges from the sequential criterion for limits, and it follows from the comparison test that $\sum f(n)$ converges. 
\end{proof}
\end{document}